\newpage
\section{Methodology}
\label{sec:method}

This section describes the architecture and training methodology of the
decoder-only Transformer. The model follows the GPT design philosophy
\parencite{radford2019language}: a stack of Pre-LayerNorm blocks, each
containing causal multi-head self-attention followed by a position-wise
feed-forward network, with residual connections throughout and weight
tying between the embedding and language-model head. Every component
below was written from raw PyTorch primitives: \texttt{nn.Linear},
\texttt{nn.LayerNorm}, \texttt{torch.matmul}, and \texttt{F.softmax},with
no high-level Transformer modules. The full implementation is roughly
$250$ lines of code across \texttt{model.py} and the supporting training
scripts.

\subsection{Decoder-only Transformer}
\label{subsec:transformer}

\subsubsection{Embedding Layer}
The embedding layer combines two components. First, a \emph{learnable
token embedding} matrix $E \in \mathbb{R}^{V \times d}$ maps each token id
in the vocabulary of size $V$ to a $d$-dimensional vector. Second, a
\emph{sinusoidal positional encoding} $P \in \mathbb{R}^{T \times d}$ is
added to inject order information:
\begin{equation}
P_{(t, 2i)}   = \sin\!\left(\frac{t}{10000^{2i/d}}\right), \qquad
P_{(t, 2i+1)} = \cos\!\left(\frac{t}{10000^{2i/d}}\right).
\label{eq:pe}
\end{equation}
Sinusoidal encodings were chosen over learned positional embeddings for
two practical reasons: they extrapolate gracefully to sequence lengths
longer than those seen during training, and they add no trainable
parameters \parencite{vaswani2017attention}. For a character-level model
on a small corpus, the parameter saving is negligible, but the
extrapolation property is worth preserving for future scaling. The final
input to the first block is $X = E[\text{ids}] + P_{:T} \in
\mathbb{R}^{B \times T \times d}$.

\subsubsection{Causal Multi-Head Self-Attention}
The core of the architecture is scaled dot-product attention, computed
manually without high-level wrappers:
\begin{equation}
\operatorname{Attention}(Q, K, V) =
\operatorname{softmax}\!\left(\frac{QK^\top}{\sqrt{d_k}} + M\right) V,
\label{eq:attention}
\end{equation}
where $M \in \{0, -\infty\}^{T \times T}$ is a causal lower-triangular
mask with zeros on and below the diagonal and $-\infty$ above. Adding
$-\infty$ before the softmax zeros out attention to future tokens,
enforcing the autoregressive property. Concretely, the mask is
constructed once as
\texttt{torch.tril(torch.ones(T, T))}, reshaped to broadcast over batch
and head dimensions, and registered as a module buffer so it moves with
the model to GPU automatically without being treated as a trainable
parameter.

\noindent For multi-head attention, $Q$, $K$, $V$ are projected $h$ times into
subspaces of dimension $d_h = d/h$, attention is computed independently
per head, and the outputs are concatenated and linearly projected:
\begin{equation}
\operatorname{MHA}(X) =
\bigl[\operatorname{head}_1 \| \cdots \| \operatorname{head}_h\bigr] W^O,
\label{eq:mha}
\end{equation}
where $\operatorname{head}_i = \operatorname{Attention}(XW^Q_i, XW^K_i, XW^V_i)$.
In our implementation, the $h$ heads are computed in parallel by a single
fused $\operatorname{Linear}(d, 3d)$ projection that produces $Q$, $K$,
and $V$ together, followed by a reshape from $(B, T, d)$ to
$(B, h, T, d_h)$. This is the same trick used in GPT-2: it is measurably
faster than three separate projections because the underlying matrix
multiply is batched, and it keeps the three weight matrices in adjacent
memory.

\subsubsection{Pre-LayerNorm Transformer Block}
Each block applies the following two sub-layers, both wrapped in a
residual connection:
\begin{align}
X' &= X + \operatorname{MHA}\!\bigl(\operatorname{LN}(X)\bigr),
\label{eq:block1}\\
X'' &= X' + \operatorname{FFN}\!\bigl(\operatorname{LN}(X')\bigr),
\label{eq:block2}
\end{align}
where $\operatorname{LN}$ is LayerNorm and the FFN is a two-layer MLP with
a GELU activation:
\begin{equation}
\operatorname{FFN}(x) = W_2 \cdot \operatorname{GELU}(W_1 x + b_1) + b_2,
\label{eq:ffn}
\end{equation}
with hidden dimension $4d$. Pre-LN (applying LayerNorm \emph{before} the
sub-layer) is preferred over Post-LN because it produces smoother
gradients, is more stable at depth, and does not require learning-rate
warmup on the residual path \parencite{xiong2020layer}. We found this
choice to matter in practice: our validation curve remains smooth from
the very first evaluation point (loss $2.47$ at step $250$) with no
early-training spikes, despite a relatively aggressive peak learning
rate of $3 \times 10^{-4}$.

\noindent The full model is: embedding $\rightarrow$ $N$ stacked blocks $\rightarrow$
final LayerNorm $\rightarrow$ linear language-model head. Following
\textcite{radford2019language}, we \emph{tie} the LM head weights to the
token embedding matrix, which reduces the parameter count and acts as a
mild regulariser by forcing the input and output embeddings to occupy the
same subspace.

\subsubsection{Optimisation and Mixed Precision}
Training uses mixed precision via \texttt{torch.autocast} with a
\texttt{GradScaler}. When running under fp16, the scaler dynamically
inflates the loss before the backward pass to prevent gradient underflow
in the lower-precision format, then unscales the gradients before
clipping. Under bf16 the scaler is disabled, because bf16 shares fp32's
exponent range and therefore does not suffer the same underflow behaviour
\parencite{micikevicius2018mixed}. In this work the model was trained
under bf16, so the scaler was active but effectively a no-op.

\noindent The learning-rate schedule is a linear warmup followed by cosine
annealing:
\begin{equation}
\eta(t) =
\begin{cases}
\eta_{\max} \cdot \dfrac{t}{T_w}, & t < T_w \\[0.6em]
\eta_{\min} + \dfrac{\eta_{\max} - \eta_{\min}}{2}
\left(1 + \cos\!\left(\pi \dfrac{t - T_w}{T - T_w}\right)\right), & t \geq T_w
\end{cases}
\label{eq:lr}
\end{equation}
where $T_w$ is the warmup length and $T$ the total number of steps. Warmup
stabilises early training when attention statistics are still noisy; the
cosine tail lets the model settle into a well-conditioned minimum without
the abrupt transition that step decay introduces.

\noindent The optimiser is AdamW \parencite{loshchilov2019decoupled}. Following the
standard GPT recipe, weight decay is applied \emph{only} to parameters
with two or more dimensions (linear and embedding weights), while 1D
parameters (LayerNorm gains and all biases) are excluded:
\begin{equation}
\theta_{t+1} = \theta_t - \eta_t
\left(\frac{\hat{m}_t}{\sqrt{\hat{v}_t} + \epsilon}
+ \lambda \theta_t \cdot \mathbb{1}[\dim(\theta) \geq 2]\right).
\label{eq:adamw}
\end{equation}
In our run this split produced $18$ decayed tensors and $34$ non-decayed
tensors. Theoretically motivated grouping of this kind matters more than
it might first appear: decaying a LayerNorm gain pulls the normalisation
scale toward zero, which directly fights the mechanism that keeps
activations well-conditioned as the network deepens.

\subsubsection{Autoregressive Generation}
At inference time, we sample from the model autoregressively:
\begin{equation}
x_{t+1} \sim p_\theta(\cdot \mid x_{\leq t}),
\label{eq:ar}
\end{equation}
where the distribution is obtained by applying softmax to the final logits
of the forward pass on the current context. Three sampling strategies are
supported:
\begin{itemize}
    \item \textbf{Greedy decoding}: $x_{t+1} = \arg\max\, p_\theta(\cdot \mid x_{\leq t})$.
          Deterministic, and prone to repetitive loops once the model
          enters a high-probability cycle.
    \item \textbf{Temperature scaling}: logits are divided by a temperature
          $T$ before softmax. $T < 1$ sharpens the distribution,
          $T > 1$ flattens it. Temperature is applied uniformly across
          the vocabulary, so it also amplifies low-probability errors
          along with the diversity of plausible continuations.
    \item \textbf{Top-$k$ sampling}: only the top-$k$ logits are retained;
          all others are set to $-\infty$ before softmax. This prevents
          sampling from the long tail of unlikely tokens while preserving
          diversity among the high-probability candidates, which is why
          it is the standard choice in most modern language models
          \parencite{holtzman2020curious}.
\end{itemize}